\section{Discussion}

\subsection{Complementarity Does Not Imply Adaptation Value}

The portfolio is complementary: three policies win somewhere by a resolved
margin, the winner map needs more than one factor, and
\NumParetoNonSingleton{} of contexts have a non-singleton Pareto set. On the
usual reading, that is the evidence for building a selector.

It is not sufficient evidence. The decision value of the entire winner map is
\NumHeadroom~s per vehicle, \NumHeadroomPct{} of the mean journey time under the
single-best fixed policy, because that policy's upside where it wins is
\NumAsymRatio$\times$ its downside where it does not. The two quantities are
computable from the same data, one is small when the other is large, and only
the second determines whether adaptation is worth attempting.

The practical implication is a cheap first step. Before choosing a model class,
compute the gap between the best fixed policy and the retrospective best, and
compare it with the noise floor of that same paired comparison. Neither requires
a model. Had we done so here, the answer would have been a \NumHeadroomPct{} gap
sitting close to the resolution of a \NumSeeds{}-seed design.

\subsection{Decision-Oriented Evaluation Changes the Interpretation of Accuracy}

Selection accuracy and decision regret rank the methods differently in this
study, and the disagreement is not marginal. The mechanistic rule is more
accurate than the fixed policy (\NumBOneTopOne{} versus \NumBZeroTopOne{}) and
more costly (\NumBOneRegret~s versus \NumBZeroRegret~s). The advantage-based
selector is less accurate than the logistic baseline and cheaper than it.

The cause is not model capacity but the training target. Methods trained to name
the winner optimise a loss that treats every context alike; a context where the
margin is 20\,s contributes no more than one where it is 0.2\,s. Predicting
advantage weights each context by what is actually at stake, which under an
\NumAsymRatio$\times$ asymmetry is the difference between a selector that helps
and one that hurts. This is the predict-then-optimise
argument~\cite{elmachtoub2022spo}, appearing here not as a refinement but as a
sign change.

For anyone reporting such a system: report decision regret, and do not present
selection accuracy as a proxy for it.

\subsection{A Strong Fixed Policy Can Be a Rational Default}

In this benchmark the load-balancing policy is optimal in
\NumSBSOptimalPct{} of contexts and loses only \NumAsymLoss~s on average where
it is not. Against that, the distance-minimising policy costs
\NumFixedPenAbsPOne~s per vehicle more --- \NumFixedPenPctAbsPOne{} of the
single-best policy's \NumSBSMeanCost~s mean journey time.

Two decisions therefore sit three orders of magnitude apart: which fixed policy
to run, and whether to vary it by context. A benchmark that reports only the
second makes the first invisible.

\subsection{Selective Adaptation and the Limits of Feature-Space Support}

The selective rule prevented the observed harms: on each of the
\NumOODNHarms{} harmful shift splits it recovered the fixed policy's regret
exactly. It did so by declining to act --- the conformal interval is wider than
the effect --- and the price is the in-distribution gain the unguarded selector
achieves.

Its limits are equally concrete. The conformal gate is calibrated under
exchangeability and so is an in-distribution instrument; the support gate is a
feature-space extrapolation detector; neither detects concept shift. Split O8
shows what that means in practice: a disruption moved to a corridor carrying
none in training left every descriptor unchanged, the support gate flagged
\NumOODEightNorthSupport{} of those contexts, and the selector acted unguarded.
It happened to do well. A support gate cannot detect a shift its feature space
does not represent, and reporting the case is more useful than asserting the
principle.

\subsection{Implications for Routing-Policy Deployment}

We state what this benchmark demonstrates, not what an operator should do.

The benchmark shows that in this environment the decision carrying the
substantial value is the choice of fixed policy, and that the incremental value
of context-dependent selection is small relative to it and close to the
resolution of the experiment. It shows that if such a selector is built, the
training target determines whether it helps, and that a support-and-confidence
gate removed the harms observed under shift while also removing the gains. And
it shows that the criterion matters as much as the policy: the policy
minimising stopped delay is not the policy minimising journey time in
\NumBestCThreeDiffersPct{} of contexts, which is a different argmin rather than
a preference weighting.

Whether these carry to a real network is not established by anything here.
